\section{Evaluation}\label{sec:eval}

% Section opening paragraph: what we're answering and the table of contents.
% RQ1: Does psi+R joint optimization outperform decoupled baselines
%      in normal operation when device heterogeneity is meaningful?
% RQ2: Does psi+R survive a mid-stream worker failure on edge?
% RQ3: Does latency-optimal placement (alpha=0) dominate throughput-
%      optimal placement (alpha=1) across the design space we can
%      vary at runtime (topology, chain mode, chain length, C)?
% RQ4: What's the cost of asynchronous chain forwarding's failure
%      attribution fallback, relative to its throughput gain?

\subsection{Experimental setup}\label{sec:eval:setup}
% - Fleet table (cf REPORT.md §1.1): on-1..on-6 + ao-1 + ao-2 + ax-1
% - Model: OPT-350M (24 layers, 1024 hidden, fp16). State that
%   OPT-1.3B's monolithic weight format made it inaccessible to
%   this fleet (Negative result §X) — would need sharded loader
%   work that's orthogonal to RADP's contribution.
% - Workload: prompt = "fox jumps over the lazy dog. Once upon a
%   time", greedy decoding (temperature=0), max_tokens varies by
%   experiment.
% - Tools: experiments/run_a3_remote.py, measure_concurrent.py,
%   run_phase3_recovery.py (cf REPORT.md §1.3).
% - Critical fixes applied: 934ea27 (OPT project_in order),
%   246a02b (safetensors prefix mismatch), 382739b (profiler
%   tokenizer + CUDA timing). All measurements post-fix.

\subsection{Profiler sanity (per-layer compute landscape)}\label{sec:eval:profiling}
% - Per-layer compute (ms) on each tier post-D2.2 fix
% - AGX Orin MAXN is 16x faster than Nano CUDA at seq=64 decode
% - Without the MAXN power-mode lock the gap collapses, hiding
%   in the noise — practical reproducibility note for systems
%   reviewers

\subsection{Headline: $\psi$+$R$ joint optimization beats decoupled baselines in normal mode}\label{sec:eval:a3b}
% From REPORT.md §4.1 (OPT-350M 3-tier, A3b' N=3):
% - TBT p50 -6.5%, throughput +8.3% vs greedy, n=300 samples per cell
% - Cost-model tie at 45.3 ms max_stage_time hides the live gap
%   that pipeline transition overhead exposes (REPORT.md §4.2)
% - Bar chart figure: TBT p50/p95 + throughput for greedy / uniform
%   / jupiter_dp / ours

\subsection{Fault recovery survives where baselines abort}\label{sec:eval:recovery}
% Two layers:
% (a) Star-topology recovery (EXP-D2.1, A2 N=5): 0/300 token loss,
%     recovery step mean 617 ms, p95 670 ms. 100 ms tight spread.
% (b) Chain-topology recovery (EXP-D3 Phase 3 + 3.2):
%     - sync chain, 4-stage, on-6 SIGKILL mid-stream: 60/60
%       coherent tokens, trailer attribution log fires correctly,
%       "chain failure for already-dead on-6[10..14]; finalising
%       recovery (rewire + replay)"
%     - async chain, same setup: 60/60 (post-TimeoutError fix),
%       but attribution falls back to head_stage when heartbeat
%       loses the race; the substitution still produces a feasible
%       plan so accuracy is preserved
%     - Mirror cache verification: 8-token Generate => exactly 16
%       mirror pushes (8 steps * 2 non-first stages), arithmetic
%       confirms the data substrate works as designed

\subsection{Headline: Latency placement strictly dominates throughput placement across 28 measurement points}\label{sec:eval:matrix}
% This is the new central finding from EXP-D3 Phase F + F.2.
% Big 4x4x2 table:
%             T+sync   T+async   L+sync   L+async
% 3-stage:
%   C=1
%   C=4
%   C=16
% 4-stage:
%   C=1
%   C=4
%   C=8
%   C=16
% Interpretation:
% (a) L > T at every one of the 28 points (highlight visually)
% (b) async > sync mostly (esp at C=16, +47% T, +17-27% L)
% (c) chain-length-independent async gain (3-stage vs 4-stage
%     deltas)
% (d) Best operating point: L+async @ C=16 = 41.7 tok/s vs worst
%     (T+sync @ C=16) = 24.5 tok/s, +71% range covered by the
%     design space.

\subsection{Pipeline parallelism vs chain length}\label{sec:eval:depth}
% 3-stage vs 4-stage L+async at matched C:
% C=1: 9.7 vs 9.3 (-4%, extra hop costs a hair on cold pipeline)
% C=4: 33.7 vs 23.9 (-29%, network/compute ratio collapses when
%      each Nano gets only 1 layer)
% C=16: 40.3 vs 41.7 (+3%, saturated)
% Lesson: stage count is not a free knob. RADP-Latency DP picks
% the stage count cost-aware; we forced 4-stage via fleet
% composition to expose this — the DP itself wouldn't have.

\subsection{Algorithmic predictions in synthetic sweep}\label{sec:eval:synth}
% From REPORT.md §2: confirms the DP's polynomial runtime and the
% specific heterogeneity ratios where greedy's rounding error
% hands us a 14-22% speedup. Not the paper's main result but
% reassures reviewers about the algorithm in isolation.

\subsection{Sensitivity studies}\label{sec:eval:sens}
% - activation_bytes calibration (REPORT.md §6): 500x overestimate
%   in the legacy default flipped the DP's choice; once corrected
%   the predictions match the live measurement
% - eager vs lazy backup (PHASES.md A5): trade-off table.
%   We default eager; lazy gives more layers to the primary at
%   the cost of slower recovery and possible token loss
% - chain_mode and optimization_mode are the headline knobs

\todo{Each subsection deserves a figure or table. Headline figures:
(i) the 28-point matrix in §6.4, (ii) recovery wall-clock CDF in
§6.3, (iii) compute-tier per-layer ms in §6.2.}
